\section{The tower and the main statements}\label{sec:statements}

Fix a field $k$ of characteristic three, an element $a\in k^*$ and a
transcendental element $t$. All splitting fields are taken in a fixed
algebraic closure of $k(t)$. Set
\[
 f(X)=X^3+aX^2,\qquad
 L_n=k(t)\bigl(f^{-n}(t)\bigr),\qquad
 G_n=\Gal(L_n/k(t)),\quad n\geq1.
\]
The notation $f^{\circ n}$ denotes composition, and $f^{-n}(t)$ is the
set of roots of $f^{\circ n}(X)-t$. The field $L_n$ is the splitting
field, not merely the field generated by one chosen root. The
polynomials are separable and irreducible; this will also follow
directly in Section~\ref{sec:normal}.

The root $t$, its three preimages, and their successive preimages form
the regular rooted ternary tree $T_\infty$. Its truncation at height
$n$ is $T_n$. A labelling assigns $1,2,3$ to the children of every
vertex; compatibility means that the labellings at finite heights are
restrictions of this one labelling. Put $W_n=\Aut(T_n)$. If
$g=(g_1,g_2,g_3)\sigma\in W_n$ and $n\geq2$, define
\[
 \sgn_2(g)=\sgn(\sigma)
             \prod_{i=1}^3\sgn(g_i|_{T_1}).
\]
Here the signs refer to permutations of three children. Following the
classical definition in \cite[Section~2]{benedetto2017cubic} and
\cite[Definition~2.2]{ejder2022arithmetic}, set
\begin{equation}\label{eq:Edefinition}
 E_1=\sym_3,\qquad
 E_n=(E_{n-1}^{\,3}\rtimes\sym_3)\cap\ker\sgn_2\quad(n\geq2).
\end{equation}
Equivalently, at every vertex having two levels below it, its local sign
equals the product of the local signs at its three children. The
restriction maps define $E=\varprojlim_n E_n$. Equality with these
groups below is equality in the labelled tree action, not just an
abstract isomorphism or a condition on the sign of all leaves.

Choose an algebraic closure $\kk$ of $k$ in the ambient algebraic
closure and put $\Lbar_n=\kk L_n$. Its group over $\kk(t)$ is the
geometric group. All local assertions in the next two geometric
statements concern $\Lbar_n/\kk(t)$.

\begin{theorem}[Arithmetic and geometric tower]\label{thm:tower}
For every $k$ and $a$ as above there is one compatible labelling such
that, for all $n\geq1$,
\begin{equation}\label{eq:grouporder}
 \Gal(L_n/k(t))=\Gal(\Lbar_n/\kk(t))=E_n,
 \qquad |E_n|=2^{3^{n-1}}3^{(3^n-1)/2}.
\end{equation}
The extension $L_n/k$ is regular; in particular its algebraic constant
field is exactly $k$. The arithmetic and geometric inverse-limit
actions both equal $E$ in the same labelling.
\end{theorem}

For the radical description, let $\mathcal A_n=f^{-n}(t)$ and, for
each $\alpha\in\mathcal A_n$, choose
$\beta_\alpha$ with $\beta_\alpha^2=\alpha/a^3$. Define
\[
 \Mbar_n=\Lbar_n(\beta_\alpha:\alpha\in\mathcal A_n),
 \qquad \wpmap(u)=u^3-u.
\]
The choice of each square-root sign does not affect the fields or the
ranks in the following statement. A \emph{sibling triple} is the set
of three roots $\alpha$ with the same parent $f(\alpha)$.

\begin{theorem}[Global and local radical ranks]\label{thm:ranks}
For every $n\geq1$, the relation space among the square classes
$\alpha/a^3$ in $\Lbar_n^*/\Lbar_n^{*2}$ is generated by the
indicator vectors of the sibling triples. Consequently
\begin{equation}\label{eq:radicaldegrees}
 [\Mbar_n:\Lbar_n]=2^{2\cdot3^{n-1}},\qquad
 [\Lbar_{n+1}:\Mbar_n]=3^{3^n}.
\end{equation}
The $3^n$ classes $[\beta_\alpha]$ in
$\Mbar_n/\wpmap(\Mbar_n)$ are linearly independent over $\F_3$.
The extension $\Mbar_n/\Lbar_n$ splits completely at every place
above infinity. At every corresponding completion of $\Mbar_n$, the
images of those same Artin--Schreier classes span exactly one
dimension over $\F_3$.
\end{theorem}

Let $\overline C_n$ be the smooth projective curve over $\kk$ with
function field $\Lbar_n$. All valuations in the next theorem are
normalized to have value group $\mathbb Z$ in the field upstairs; a
different exponent is measured in that normalization.

\begin{theorem}[Geometric branch data and genus]\label{thm:genus}
The cover $\overline C_n\to\PP^1_{\kk}$ has branch set exactly
$\{0,\infty\}$. At every point above zero, its ramification index
and different exponent are
\begin{equation}\label{eq:zero-data}
 e_{0,n}=2^n,\qquad d_{0,n}=2^n-1.
\end{equation}
At every point above infinity they are
\begin{equation}\label{eq:infinity-data}
 e_{\infty,n}=2\cdot3^n,\qquad d_{\infty,n}=3^{n+1}-2.
\end{equation}
Its genus is
\begin{equation}\label{eq:genus}
 g_n=1+\frac{|E_n|}{2}
       \left(\frac12-\frac1{2^n}-\frac1{3^n}\right).
\end{equation}
In particular, $g_1=0$ and $g_2=46$.
\end{theorem}

Theorem~\ref{thm:tower} is over an arbitrary, possibly imperfect,
constant field. Theorems~\ref{thm:ranks} and \ref{thm:genus} are
geometric statements after adjoining $\kk$; they do not assert that
all branch points over the original $k$ have residue degree one.
The height $n$ is neither a forward-period clock nor a finite-field
extension degree. No surjectivity statement for a specialization
$t=t_0\in k$ is included.

The proof proceeds simultaneously in the group and the infinity
index. The field-independent signature bound is proved first. Assuming
the result at height $n$, the zero-place valuations determine the next
square-class rank, while the already known infinity index gives pole
order one for the square roots. Character projections prove their
global independence, and a separate leading-coefficient calculation
proves local rank one. The two conclusions then establish the next
group and the next infinity index. This order avoids using the
unproved group at height $n+1$ in the rank calculation.
